% ATTEMPT_STATUS: unresolved
% TOP_PROBLEM: {"problemId": "problem.colliot-thelene-conjecture-brauer-manin-obstruction-controls-rational-points-on-rationally-connected-varieties", "canonicalTitle": "Colliot-Thelene Conjecture: Brauer-Manin Obstruction Controls Rational Points on Rationally Connected Varieties", "releaseRank": 151, "primaryDomain": "number_theory_arithmetic_geometry", "primaryDomainLabel": "Number theory & arithmetic geometry", "displayStatus": "open_with_solved_subcases", "releaseStatus": "open_with_solved_subcases"}
% !TeX program = lualatex
% Complete problem section copied verbatim from research_batch_151-153.tex.
\documentclass[11pt]{article}
\usepackage[margin=25mm]{geometry}
\usepackage{fontspec}
\setmainfont{Latin Modern Roman}
\usepackage{amsmath,amssymb,mathtools}
\usepackage{unicode-math}
\setmathfont{Latin Modern Math}
\usepackage{enumitem,needspace}
\usepackage{xurl,hyperref}
\hypersetup{colorlinks=true,urlcolor=blue}
\setcounter{secnumdepth}{1}
\setlength{\parindent}{0pt}
\setlength{\parskip}{5pt}
\setlength{\emergencystretch}{3em}
\setlist[itemize]{leftmargin=3em,itemsep=5pt,topsep=4pt,parsep=0pt}
\allowdisplaybreaks
\newcommand{\N}{\mathbb{N}}
\newcommand{\Z}{\mathbb{Z}}
\newcommand{\Q}{\mathbb{Q}}
\newcommand{\R}{\mathbb{R}}
\newcommand{\C}{\mathbb{C}}
\newcommand{\F}{\mathbb{F}}
\newcommand{\A}{\mathbb{A}}
\newcommand{\PP}{\mathbb P}
\newcommand{\E}{\mathbb{E}}
\newcommand{\eps}{\varepsilon}
\newcommand{\dd}{\,\mathrm d}
\newcommand{\sref}[2]{\hyperlink{source:#1:#2}{[S#2]}}
\newcommand{\eref}[2]{\hyperlink{extra:#1:#2}{[E#2]}}
% Turn the next switch off for a shorter reading copy. The quotations preserve
% every exactTarget field, including qualifications omitted from a short summary.
\newif\ifshowcatalogue
\showcataloguetrue
\newcommand{\cataloguescope}[1]{\ifshowcatalogue
 \paragraph{Exact scope in the supplied catalogue.}
 \begin{quote}\small #1\end{quote}\fi}
\begin{document}
\setcounter{section}{150}
% ================================================================
% problemId: problem.colliot-thelene-conjecture-brauer-manin-obstruction-controls-rational-points-on-rationally-connected-varieties
% source releaseRank: 151; source releaseStatus: open_with_solved_subcases
% exactTarget SHA-256: 5d419b183433689d7c93892513ac35d31e51d57820703ac1acc63921dc0158d1
\clearpage
\section{\texorpdfstring{Colliot-Thelene Conjecture: Brauer-Manin Obstruction Controls Rational Points on Rationally Connected Varieties}{Colliot-Thelene Conjecture: Brauer-Manin Obstruction Controls Rational Points on Rationally Connected Varieties}}\label{151}
\subsection{Definitions and mathematical statement}
For a number field $k$, write $\A_k$ for its ring of adeles and $\operatorname{Br}(X)=H^2_{\mathrm{et}}(X,\mathbf G_m)_{\mathrm{tors}}$. Evaluation of a Brauer class at local points, followed by the local invariant maps, defines
\[
 X(\A_k)^{\operatorname{Br}}=\{(x_v):\ \sum_v\operatorname{inv}_v(b(x_v))=0\text{ for every }b\in\operatorname{Br}(X)\}.
\]
A geometrically rationally connected variety has rational curves joining two general geometric points. For every smooth proper geometrically irreducible such $X$, the assertion is
$\overline{X(k)}=X(\A_k)^{\operatorname{Br}}$ in the adelic topology. Thus a compatible adelic point should be approximable by rational points at any prescribed finite collection of places. The precise adelic convention, including archimedean factors, is that of the cited formulation.
\par\smallskip\textit{Formulation and concept references:} \sref{151}{1}.
\cataloguescope{Let X be a smooth, proper, geometrically irreducible variety over a number field k whose base change to an algebraic closure is rationally connected. Prove that X(k) is dense in the Brauer-Manin set X(A\_k)\textasciicircum{}\{Br(X)\}; in particular the Brauer-Manin obstruction is the only obstruction to the Hasse principle and weak approximation for such varieties.}
\subsection{Short English statement}
For a rationally connected variety over a number field, do the Brauer compatibility conditions account for every failure of rational points to exist or approximate local solutions?
% BEGIN RESEARCH ADDITION ATTEMPT-151
\subsection{Research attempt}

\paragraph{Scope check and current frontier (24 September 2026).}
The target includes approximation at the archimedean places in their ordinary
local topology.  Conjecture 2.3 of \sref{151}{1} uses $X(\A_k)$, not the modified
space $X(\A_k)_\bullet$ introduced immediately before it.  We retain this stronger
approximation requirement.  The quotient
$B_X=\operatorname{Br}(X)/\operatorname{Br}_0(X)$ is finite for the varieties in
this entry; here $\operatorname{Br}_0$ denotes the image of
$\operatorname{Br}(k)$.  See \sref{151}{1}, Remarks 2.4(i)--(ii), and
\eref{151}{1}, Lemmas 1.1 and 1.3.
This makes the obstruction finite, but does not make solving it by rational
points automatic.

The established fibration theorem of Harpaz--Wittenberg propagates the desired
density from rationally connected fibres over projective space, conditional on
their conjecture about locally split polynomial values; its unconditional cases
require additional degeneration hypotheses.  The zero-cycle theorem in the
same paper is not a rational-point theorem
(\eref{151}{2}, Theorems 1.3, 1.5--1.6 and Corollary 9.25).
Recent connected-group descent results of Nguyen Manh Linh give further
propagation theorems, not a proof for every rationally connected variety
(\eref{151}{3}, Theorem A and \S3.4).
Tian's 2026 manuscript gives an alternative descent argument for certain
connected-group torsors; its consulted version is recorded as a preprint, not
as an independently audited solution of this entry (\eref{151}{4}, abstract
and Theorem 1.4).

\paragraph{Attack map.}
Three routes were considered.  First, deformation of rational curves can move
points after a suitable curve and its arithmetic marking have been constructed;
the missing assertion is the existence of a \emph{$k$-defined} curve with the
required simultaneous local markings.  Geometric rational connectedness alone
does not supply that assertion.  Second, descent could replace $X$ by twists
of torsors, but a proof still needs the desired density on those twists; using
the universal conjecture there would be circular.  Third, a fibration could
move an adelic point onto a rational fibre and then solve the fibre.  Here one
must simultaneously control local solubility and \emph{all} fibre Brauer
classes, including classes not restricted from $X$; this is the issue addressed
by \eref{151}{2}, \S9.

The route pursued below is the third, with a finite Fourier calculation added
to the arithmetic fibration problem.  The proposed benefit is that several
places with individually small evaluation images might together correct the
entire remaining Brauer defect.  The calculation is proved below; the existence
of suitable places and parameters is deliberately left as an arithmetic
hypothesis.  This is an attempted combination of methods, not a claim that
the finite-group calculation is new in harmonic analysis.

\paragraph{Attempt: separating the inherited obstruction.}
For this route only, suppose that a smooth proper geometrically integral variety $Y/k$ has a dominant
morphism $f:Y\to\PP^1_k$ with geometrically rationally connected generic fibre.
Assume also that $\operatorname{Br}(Y)/\operatorname{Br}_0(Y)$ is finite; this
holds, in particular, when $Y$ is itself geometrically rationally connected.
We do not assert that every $X$ in the problem has such a presentation.
Fix $a=(a_v)\in Y(\A_k)^{\operatorname{Br}}$ and a basic open neighbourhood of
$a$, prescribed at a finite set $S_0$ of places.  Suppose, as a separate first
step, that a smooth geometrically rationally connected fibre $Z=Y_t$ over
$t\in\PP^1(k)$, and an adelic point
$z^0=(z_v^0)\in Z(\A_k)$ have been produced, with $z_v^0$ in the required
neighbourhoods for $v\in S_0$.

Put
\[
 B=\operatorname{Br}(Z)/\operatorname{Br}_0(Z),\qquad
 C=\operatorname{im}\bigl(\operatorname{Br}(Y)\longrightarrow B\bigr),
 \qquad B^\vee=\operatorname{Hom}(B,\Q/\Z).
\]
These are finite groups.  The total defect
\begin{equation}\label{ra151:defect}
 e([b])=\sum_v\operatorname{inv}_v b(z_v^0),\qquad [b]\in B,
\end{equation}
is a well-defined element of $B^\vee$: changing $b$ by a class from
$\operatorname{Br}(k)$ changes the sum by zero, by global reciprocity.  The
sum is finite.  Indeed, finitely many torsion representatives of $B$ extend
over a proper integral model away from finitely many places; their evaluations
there lie in $\operatorname{Br}(\mathcal O_v)=0$.  These are the usual
spreading-out and properness facts underlying Brauer--Manin evaluation;
compare \eref{151}{1}, \S1.  No convergence of an infinite conditionally
convergent series is involved.

Enlarge $S_0$ to a finite set $S$ containing the archimedean places and the
supports of chosen representatives of
$\operatorname{Br}(Y)/\operatorname{Br}_0(Y)$.  If $z_v^0$ is sufficiently close
to $a_v$ at $S$, local constancy of evaluation gives
\[
 e|_C=0.
\]
To check this, take a representative $\beta\in\operatorname{Br}(Y)$ from the
chosen finite collection.  Its evaluations at $a_v$ and at $z_v^0$ agree on
$S$, and both vanish off $S$.  Their global sums therefore agree and are zero.
Taking integer combinations and using reciprocity for constant classes proves
the assertion for every element of $C$.  Thus the remaining defect belongs to
\[
 H=\operatorname{Ann}(C)=\{h\in B^\vee:h(c)=0\text{ for all }c\in C\}.
\]
It is important to work in $H$, rather than trying to mix on all of $B^\vee$.

For $v\notin S$ define the \emph{difference} evaluation map
\begin{equation}\label{ra151:local-difference}
 d_v:Z(k_v)\longrightarrow H,\qquad
 d_v(z)([b])=\operatorname{inv}_v\bigl(b(z)-b(z_v^0)\bigr),
 \qquad A_v=d_v(Z(k_v)).
\end{equation}
Constants cancel locally, so this map is independent of the lift of $[b]$.
It takes values in $H$ because the inherited classes have constant evaluation
on $Y(k_v)$ for $v\notin S$.  Every $A_v$ is a finite set containing zero.
For a finite set $T$ of places disjoint from $S$, keeping $z^0$ unchanged away
from $T$ yields the exact equivalence
\begin{equation}\label{ra151:sumset}
 \text{the defect can be corrected using $T$}
 \quad\Longleftrightarrow\quad
 -e\in\sum_{v\in T}A_v.
\end{equation}
Here the right side is a sum of sets, with one choice from each set; it is
\emph{not} the subgroup generated by their union.  This distinction will be
an essential stress test.

\paragraph{Finite-character repair lemma (proved here).}
Let $H$ be a finite abelian group, let $A_1,\ldots,A_m\subseteq H$ be nonempty,
and choose probability weights $\mu_i$ supported on $A_i$.  For a complex
character $\chi\in\widehat H=\operatorname{Hom}(H,\{z\in\C:|z|=1\})$, write
\[
 \widehat\mu_i(\chi)=\sum_{h\in H}\mu_i(h)\chi(h).
\]
If
\begin{equation}\label{ra151:fourier-condition}
 \mathcal R:=\sum_{\substack{\chi\in\widehat H\\\chi\ne1}}
                  \prod_{i=1}^m|\widehat\mu_i(\chi)|<1,
\end{equation}
then $A_1+\cdots+A_m=H$.  More precisely, every $h\in H$ has probability at
least $(1-\mathcal R)/|H|$ under the convolution $\mu_1*\cdots*\mu_m$.

\emph{Proof.}
Character orthogonality gives, for independent $H$-valued random variables
$D_i$ of laws $\mu_i$,
\[
 \Pr\left(\sum_iD_i=h\right)
 =\frac1{|H|}\sum_{\chi\in\widehat H}
       \overline{\chi(h)}\prod_i\widehat\mu_i(\chi).
\]
All sums and products are finite.  The trivial character contributes $1/|H|$;
the absolute value of the remaining contribution is at most
$\mathcal R/|H|$.  Positivity proves the assertion.  Independence here is
merely a convenient way of enumerating the Cartesian product of local choices;
it is not an unproved independence assertion about primes.\hfill$\square$

A useful sufficient hypothesis is the following quantitative one.  Suppose
$0<\eta<1$ and every nontrivial $\chi$ has at least $r$ indices $i$ for which
$|\widehat\mu_i(\chi)|\le1-\eta$.  Then
\begin{equation}\label{ra151:quantitative}
 (|H|-1)(1-\eta)^r<1
 \quad\Longrightarrow\quad A_1+\cdots+A_m=H.
\end{equation}
The indices may depend on $\chi$.  The proof uses
$|\widehat\mu_i(\chi)|\le1$ at all other indices and sums the resulting bound.
If $H=0$, there is no defect left and the character sum is empty.

\paragraph{Conditional application to a fibre.}
In \eqref{ra151:local-difference}, choose finitely many local points realizing
weights $\mu_v$ on $A_v$, $v\in T$.  If
\eqref{ra151:fourier-condition} holds on $H$, the lemma and
\eqref{ra151:sumset} produce $z\in Z(\A_k)^{\operatorname{Br}}$ without changing
any prescribed component at $S$.  If, \emph{in addition}, rational points of
this fibre are dense in its Brauer--Manin set in the full local topology, a
point of $Z(k)\subset Y(k)$ lies in the original neighbourhood of $a$.
There is no assumption that the restriction $\operatorname{Br}(Y)\to
\operatorname{Br}(Z)$ is surjective.

Consequently, the following precise conditional propagation statement has
been obtained.  If for every $a$, every prescribed finite neighbourhood and
every resulting enlargement $S$, one can select a smooth geometrically
rationally connected fibre over a rational parameter, local points at \emph{all} places, and a finite set $T\cap S=\varnothing$
satisfying the hypotheses above, and the selected fibres satisfy the target
density statement, then $Y(k)$ is dense in $Y(\A_k)^{\operatorname{Br}}$.
The proof is the preceding two paragraphs.  Selection of $t$, existence of
$z^0$, the mixing inequality, and the fibre theorem are separate hypotheses;
none has been derived for an arbitrary fibration.

\paragraph{Stress test: where the new route breaks.}
\emph{Generation is not correction.}  In $H=\Z/3\Z$, the set $A=\{0,1\}$
generates $H$ but does not contain $2$.  A single place with this image cannot
correct a defect equal to $1$.  One cannot sample this same place twice and
pretend to have two independent adelic coordinates.  In contrast, for
$H=(\Z/2\Z)^2$, the uniform measures on $\{0,(1,0)\}$ and $\{0,(0,1)\}$
have uniform convolution on $H$.  These are exact finite examples, not
numerical evidence for an arithmetic assertion.

\emph{Good primes do not provide unlimited repair.}  For a \emph{fixed}
proper fibre $Z$, all the difference maps $d_v$ vanish away from a finite set,
by the same integral-model argument used for \eqref{ra151:defect}.  Adding
arbitrarily many good primes merely adds factors
$|\widehat\mu_v(\chi)|=1$.  Thus the attractive shortcut ``use enough
Chebotarev primes and mix'' fails unless varying $t$ makes those primes genuinely
active for the fibre's extra classes.  Good reduction bounds make this
obstruction structural, not just an inconvenience in a particular example
(\eref{151}{1}, Theorem 3.1).

\emph{Congruences do not control all prime divisors.}  For instance, in a norm
fibration $N_{L/k}(u)=P(t)$, after choosing $t$ by finitely many congruences,
new prime divisors of $P(t)$ can still obstruct local solubility.  At a place
where $L\otimes_k k_v=L_w$ is an unramified field extension of residue
degree $f>1$, valuations of norms are divisible by $f$; a newly appearing valuation one is therefore an
obstruction.  The valuation statement follows from
$v_k(N_{L_w/k_v}u)=f(w/v)v_w(u)$, with normalized valuations.  Controlling
all newly appearing places is precisely the kind of simultaneous split-value
problem in \eref{151}{2}, Conjecture 9.1.  A Fourier estimate after $t$ is
chosen does not supply that missing selection theorem.  In particular, $H=0$
does not by itself imply that a locally soluble fibre exists.

\emph{A fibration is not a universal dimensional reduction.}  Taking
$Y=X\times\PP^1$ gives fibres isomorphic to $X$ and is circular.  Taking an
arbitrary pencil on $X$ need not give rationally connected fibres.  For example,
on a smooth cubic surface, a general hyperplane section $D=-K_X$ has
$2g(D)-2=(K_X+D)D=0$, hence genus one by adjunction.  This computation, using $K_X=-H$ for a cubic surface in $\PP^3$, shows
why a generic anticanonical pencil does not establish an induction even in
dimension two.  It does not assert that no useful special fibration exists.
Similarly, replacing a point by a formal zero-cycle of degree one would not
prove existence of an effective degree-one cycle: negative coefficients cannot
be discarded.

\paragraph{Outcome.}
\textbf{Outcome: Exploratory approach with unresolved gap.}

The exact finite-defect equation, the inherited-class annihilator, and a
sufficient Fourier criterion for repairing all remaining fibre classes have
been derived with complete elementary proofs.  They give a conditional
propagation argument for specifically equipped fibrations.  No new class of
varieties satisfying the conjecture is established here, and no priority claim
is made for the finite-group lemma.

To turn this route into a proof one still needs an arithmetic selection theorem
producing locally soluble rational fibres with enough \emph{distinct active}
places, or a different argument correcting the particular defect without the
stronger full-support condition.  The fibre density hypotheses must also be
proved, and a route reaching arbitrary rationally connected varieties must be
supplied.  These missing assertions are not consequences of geometric rational
connectedness or finiteness of the Brauer group alone.  No counterexample to
the exact target was obtained.
% END RESEARCH ADDITION ATTEMPT-151
\subsection{Sources}
\begin{itemize}\raggedright
\item[\textnormal{[S1]}]\hypertarget{source:151:1}{} Olivier Wittenberg, "Rational points and zero-cycles on rationally connected varieties over number fields," arXiv:1604.08543; Proc. Sympos. Pure Math. 97.2 (2018), 597-635.
\url{https://arxiv.org/pdf/1604.08543}.
{\small\textit{Catalogue formulation source.}}
% BEGIN RESEARCH ADDITION REFERENCES-151
\item[\textnormal{[E1]}]\hypertarget{extra:151:1}{} J.-L. Colliot-Th\'el\`ene and A. N. Skorobogatov, \emph{Good reduction of the Brauer--Manin obstruction}, Transactions of the American Mathematical Society 365 (2013), 579--590; Lemmas 1.1 and 1.3, \S1 and Theorem 3.1. \url{https://arxiv.org/abs/1009.0247}.
{\small\textit{Added research source; primary manuscript consulted 24 September 2026.}}
\item[\textnormal{[E2]}]\hypertarget{extra:151:2}{} Yonatan Harpaz and Olivier Wittenberg, \emph{On the fibration method for zero-cycles and rational points}, Annals of Mathematics 183 (2016), 229--295; Theorems 1.3, 1.5--1.6, Conjecture 9.1 and Corollary 9.25. \url{https://arxiv.org/abs/1409.0993}.
{\small\textit{Added research source; the rational-point theorem is used with its explicit conditional hypotheses. Consulted 24 September 2026.}}
\item[\textnormal{[E3]}]\hypertarget{extra:151:3}{} Nguyen Manh Linh, \emph{On the descent conjecture for rational points and zero-cycles}, Journal of the Institute of Mathematics of Jussieu 25(1) (2026), 515--563, published online 19 November 2025; Theorem A and \S3.4. \url{https://doi.org/10.1017/S1474748025101400}. Author manuscript: \url{https://arxiv.org/html/2305.13228v3}.
{\small\textit{Added research source; publication metadata and author manuscript consulted 24 September 2026.}}
\item[\textnormal{[E4]}]\hypertarget{extra:151:4}{} Yisheng Tian, \emph{A Simpler Approach to a Descent Conjecture of Wittenberg}, arXiv:2604.08146 (2026), consulted HTML version, abstract and Theorem 1.4. \url{https://arxiv.org/html/2604.08146}.
{\small\textit{Added frontier reference; preprint, not an independent proof endorsement. The consulted abstract restricts the result to certain connected-group torsors. Consulted 24 September 2026.}}
% END RESEARCH ADDITION REFERENCES-151
\end{itemize}

\end{document}
